% from ega4-17-fr.tex
\oldpage[IV-4]{76}

  \item[b)] Il existe un schéma affine noethérien $S_0=\Spec(A_0)$, un morphisme de type fini $f_0:X_0\to S_0$, et un morphisme $S\to S_0$ tels que le $S$-préschéma $X_0\otimes_{S_0}S$ soit $S$-isomorphe à $X$ et que $f_0$ soit lisse (resp. non ramifié, resp. étale).
  \item[c)] Les conditions de b) sont vérifiées, et en outre $A_0$ est une sous-$\mathbb{Z}$-algèbre de type fini de $A$, le morphisme $S\to S_0$ correspondant à l'injection canonique $A_0\to A$.
\end{enumerate}
\end{corollary}

La démonstration à partir de (17.7.8) est la même que celle de (11.2.7) à partir de (11.2.6).

\begin{proposition}[17.7.10]
\label{IV.17.7.10}
Soient $f:Y\to S$, $h:X\to S$ deux morphismes localement de présentation finie, $g:X\to Y$ un $S$-morphisme, $x$ un point de $X$, $y=g(x)$. Supposons que $h$ soit plat au point $x$ et que $f$ soit non ramifié au point $y$. Alors $f$ est étale au point $y$ et $g$ est plat au point $x$.
\end{proposition}

(On verra plus loin (18.4.9) que l'on peut en fait se dispenser de l'hypothèse que $h$ est localement de présentation finie).

La question étant locale sur $S$, $X$ et $Y$, on peut supposer que $S$, $X$ et $Y$ sont affines, $f,g,h$ des morphismes de présentation finie, $h$ plat et $f$ non ramifié. Compte tenu de (11.2.7) et (17.7.9), on peut supposer en outre que $S$, $X$ et $Y$ sont noethériens; enfin on peut se borner au cas où $S=\Spec(\sh{O}_{S,s})$ (où $s=f(y)=h(x)$). L'anneau $A=\sh{O}_{S,s}$ étant un anneau local noethérien, il existe un anneau local noethérien $B$ complet, de corps résiduel algébriquement clos et un homomorphisme local $A\to B$ faisant de $B$ un $A$-module fidèlement plat ($0_{\mathrm{III}}$, 10.3.1). Remplaçant $X$ et $Y$ par $X\times_S\Spec(B)$ et $Y\times_S\Spec(B)$, on conclut par (2.5.1) et (17.7.1) qu'on peut se borner au cas où $A$ est complet et a un corps résiduel algébriquement clos. L'hypothèse que $f$ est non ramifié au point $y$ entraîne alors (17.4.4) que $\sh{O}_{Y,y}$ est une $\sh{O}_{S,s}$-algèbre finie et un anneau local noethérien complet ($0_{\mathrm I}$, 7.4.2), et comme le corps résiduel de $\sh{O}_{S,s}$ est algébriquement clos, l'homomorphisme $\sh{O}_{S,s}\to\sh{O}_{Y,y}$ est surjectif (17.4.4). Mais, d'autre part, l'hypothèse que $h$ est plat au point $x$ entraîne que l'homomorphisme composé $\sh{O}_{S,s}\to\sh{O}_{Y,y}\to\sh{O}_{X,x}$ est injectif ($0_{\mathrm I}$, 6.5.1), donc l'homomorphisme $\sh{O}_{S,s}\to\sh{O}_{Y,y}$ est bijectif, ce qui montre que $f$ est étale au point $y$ (17.6.3); en outre, $\sh{O}_{X,x}$ est un $\sh{O}_{Y,y}$-module plat, donc $g$ est plat au point $x$.

\begin{proposition}[17.7.11]
\label{IV.17.7.11}
Soient $S$ un préschéma, $X$, $Y$ deux $S$-préschémas localement de présentation finie sur $S$, $f:X\to Y$ un $S$-morphisme. Pour tout $s\in S$, on note $X_s$, $Y_s$, $f_s$ les préschémas et le morphisme déduits de $X$, $Y$, $f$ par le changement de base $\Spec(k(s))\to S$. Alors:
\begin{enumerate}
  \item[(i)] L'ensemble des $x\in X$ tels que, si $s$ est l'image de $x$ dans $S$, $f_s:X_s\to Y_s$ soit lisse (resp. non ramifié, resp. étale, resp. différentiellement lisse) au point $x$, est localement constructible.
  \item[(ii)] Supposons $f$ de présentation finie. L'ensemble des $y\in Y$ tels que, si $s$ est l'image de $y$ dans $S$, $f_s$ soit lisse (resp. non ramifié, resp. étale, resp. différentiellement lisse) en tous les points de $f_s^{-1}(y)$, est localement constructible.
  \item[(iii)] Supposons $X$ et $Y$ de présentation finie sur $S$. L'ensemble des $s\in S$ tels que $f_s$ soit lisse (resp. non ramifié, resp. étale, resp. différentiellement lisse) est localement constructible.
\end{enumerate}
\end{proposition}

Soit $E$ l'ensemble des $x\in X$ vérifiant la propriété considérée dans (i). Alors l'ensemble $F$ des $y\in Y$ vérifiant la propriété correspondante considérée dans (ii) n'est

